% Transcription — fonds Alexandre Grothendieck, shelfmark 16, batch 1 (pages 1-20).
% Established from the facsimile archives/batches/16/batch-01.pdf (a mirror of
% https://grothendieck.umontpellier.fr/16.pdf), per the skill
% transcribe-grothendieck. The batch file carries no cover sheet: PDF page N is
% archivists' page N. Checked against the pencilled numbers bottom-left — a "1"
% on the first leaf, "11", "13", "15" on the corresponding leaves and an "18",
% "19", "20" on the last three confirm the alignment end to end.
%
% Pass: Opus 5 (claude-opus-5), 2026-09-19 — first pass, unchecked against the pages by a human.
%
% What the batch is. One continuous draft in his hand, sixteen leaves (pages 2
% to 17), headed « Sur certains anneaux gradués » and paginated by him 1 to 15,
% with a "5 bis" inserted at page 7 (his §5, written after §6 had begun). Then
% three loose leaves: pages 18 and 20 are two leaves of a typewritten text of
% his own, numbered 45 and 42 in its own pagination, carrying corrections and
% insertions in his hand; page 19 is the back of such a leaf and carries a sheet
% of his own calculations on L and Λ. Page 1 is blank apart from show-through
% and gets no \page{}. Nothing in the batch is dated and no letter appears in
% these twenty pages.
%
% The run's argument, section by section (his numbering):
%
%   2 (his 1)     the data: M* a graded k-module, ℰ* a graded subring of
%                 End(M*) containing the projectors π_i; the induced grading
%                 ℰ^(ν) = Σ π_{i+ν} ℰ* π_i and the characterisation
%                 u ∈ ℰ^(ν) ⟺ u π_i = π_{i+ν} u; the relations (1.1).
%   3 (his 2)     §2, the converse: which conditions on a family (π_i) in a
%                 graded ring make (1.3) define the grading; the bigrading
%                 ℰ = ⊕ π_j ℰ π_i. §3 opens on existence and uniqueness.
%   4-5 (3-4)     §3, uniqueness of the π_i under (3.0) π_i = 0 for i ∉ [0,2n]:
%                 ℰ^0 is the commutant of the π_i, the π_i are orthogonal
%                 idempotents of the commutative ring 𝔷ℰ^0, the partial sums
%                 φ_i, ψ_{2n-i}, and the reduction to knowing ℰ^{2n-i}_i;
%                 the conditions (3.1)_i, (3.2)_i on annihilators. §4 gives
%                 sufficient conditions (4.1)-(4.4) on M*.
%   6 (his 5)     (4.4)_i ⟹ (3.1)_i and (3.2)_i, proved inside ℰ alone;
%                 §6 opens with L ∈ ℰ^(2).
%   7 (5 bis)     §5, existence: if ℰ sits in a larger graded ring ℱ where the
%                 π_i already live, (4.4)_i forces π_i ∈ ℰ, by the elements
%                 v'_i = π_{2n-i} v_i π_i and w'_i = π_i w_i π_{2n-i}.
%   8-9 (6-7)     the hard-Lefschetz condition (6.1)_i, L^{n-i} : M^i → M^{2n-i}
%                 an isomorphism; the primitive decomposition (6.4) and the
%                 projectors π_{i,α}; a Théorème listing equivalent conditions
%                 (i) to (vi), closing on (6.7), the two identities
%                 LΛ = 1 − Σ_{i≤n} π_{i,0}, ΛL = 1 − Σ_{i≥n} π_{i,0}
%                 that determine Λ.
%   10-13 (8-11)  the proof: descending induction on i, q_{i,α} = Σ_{α'≤α}π_{i,α'},
%                 the formulas (6.9), (6.10) for π_{i,α} and π_{2n-i,α},
%                 then Λ ∈ ℰ through (6.11)-(6.13); a Remarque restating the
%                 whole characterisation in L and Λ alone, (1)-(4); (6.15).
%   13-14 (11-12) §7, the model: M*_0 = Λ M_0 with ξ_0 ∈ Λ²M_0, L_0 = ξ_0 ∧ −,
%                 Λ_0 = ξ*_0 ⌟ −, and Φ_0 = Z[L_0, Λ_0] ⊂ End(M*_0); Φ_0 is
%                 generated by the φ^{β,α}_i and the multiplication table is
%                 written out.
%   14-16 (12-14) a Proposition: (i) L satisfies the Lefschetz conditions,
%                 (ii) Λ exists, (iii) there is a graded ring homomorphism
%                 φ : Φ_0 → ℰ with φ(L_0) = L, and a Corollaire on its
%                 uniqueness; then a Question, answered « Non » by n = 1,
%                 where Z[L] is only Z + ZL: L alone will not do, the π_i
%                 have to be given with it.
%   17 (his 15)   §8 opens on necessary conditions and the chain
%                 ℰ^(2n) ≠ 0 ⟺ π_{2n} ≠ 0, L^n ≠ 0, Λ^n ≠ 0.
%
% Notation fixed for the batch. His script E is the graded ring, written
% \mathcal{E}^{*}; \mathcal{E}^{(\nu)} is the homogeneous part of degree ν (he
% keeps the parentheses) and \mathcal{E}^{j}_{i} = π_j ℰ π_i the bigraded
% component, which he also writes ℰ_{2n-i,i} once, at page 4. \mathcal{E}^{0}
% without parentheses is his own shorthand for ℰ^(0) and is kept. 𝔷ℰ^0 is the
% centre, transcribed \mathfrak{z}\mathcal{E}^{0}. M^{*} is the graded module,
% P^{i} its primitive part, π_i the projector on M^i and π_{i,α} the projector
% on L^α P^{i-2α}. L is his Lefschetz operator (degree 2), Λ its inverse of
% degree −2; L_0, Λ_0, Φ_0 carry a subscript zero in §7 for the exterior-algebra
% model. v_i, w_i are the elements of (4.4); φ_i = Σ_{α≥i} π_α and
% ψ_{2n-i} = Σ_{α≤2n-i} π_α the partial sums of §3, q_{i,α} = Σ_{α'≤α} π_{i,α'}
% those of §6; φ^{β,α}_i are the generators of Φ_0 in §7 — the two φ's are his
% and are not the same letter twice by accident, the second always carrying two
% superscripts. Z is set as \mathbb{Z}.
%
% What does not carry. The leaves fall into two groups, exactly as
% references/hand.md describes: pages 2 to 8 are written fair and read almost
% end to end, whereas pages 9 to 17 are a fast palimpsest under heavy
% cancellation in which the displayed formulas carry and the connective prose
% very largely does not. The apparatus on those pages is dense in consequence,
% and several equation labels are overwritten past recovery.
%
% Revision. On 2026-09-19 the closing sentence of page 13 was re-read from the
% tiles by Opus 5 (claude-opus-5). The first pass left it largely \ill{}; it in
% fact reads « ... satisfait à la condition de (6.1)(w), et que Lambda_0 est
% l'opérateur associé défini par (6.15) », which identifies the contraction of
% the displayed formula with the operator the relations (6.15) characterise.
% The two are not the same operator (see the modernised reading), but the
% identification is what the leaf asserts, and it now stands in the record.
%
% Batch boundary. The run does not end at page 17: §8 has just been opened and
% breaks off on a list of necessary conditions, so the argument continues into
% batch 2. A \note at page 17 says so.

\documentclass[11pt,a4paper]{article}
% Le préambule commun à toutes les transcriptions du fonds Grothendieck.
%
% Il tient en une idée : **l'incertitude se compose, elle ne se résout pas.**
% Une transcription qui lisse un mot douteux a détruit la seule chose pour
% laquelle on la faisait — un lecteur doit pouvoir savoir, à chaque endroit,
% ce qui était sur le feuillet et ce qui a été ajouté. Les six macros qui
% suivent sont donc l'appareil critique, et non de la décoration.
%
% Les mêmes commandes sont reconnues par `scripts/render.mjs`, qui produit la
% vue de lecture du site. Une macro ajoutée ici doit l'être aussi là-bas,
% faute de quoi le rendu échouera bruyamment — ce qui est voulu : un rendu
% silencieusement faux est pire qu'un rendu absent.

\ProvidesPackage{grothendieck}[2026/08/08 Transcriptions du fonds Grothendieck]

\usepackage{amsmath,amssymb,amsthm}
\usepackage{mathrsfs}
\usepackage{stmaryrd}
% Les diagrammes commutatifs sont partout dans ce fonds : c'est la langue
% même de la géométrie algébrique de Grothendieck, et une transcription qui
% les rendrait en prose ne transcrirait rien.
\usepackage{tikz-cd}
% Français par défaut — c'est la langue des feuillets — mais l'anglais est
% chargé aussi : la traduction et le résumé sont des documents anglais, et la
% typographie française (espaces avant les deux-points) y serait une faute.
% Ces documents basculent par \selectlanguage{english} en tête de corps.
\usepackage[english,french]{babel}
\usepackage{fontspec}
\usepackage{geometry}
\geometry{a4paper,margin=25mm,marginparwidth=28mm}
\usepackage{marginnote}
\usepackage{xcolor}
% ulem, not soul. `soul`'s \st and \ul are built on a character-by-character
% scan that XeLaTeX's font machinery breaks: the document fails with an
% undefined control sequence rather than degrading. ulem does the same two jobs
% and is Unicode-engine safe. `normalem` stops it hijacking \emph, which the
% transcriptions use for Grothendieck's own emphasis.
\usepackage[normalem]{ulem}
\usepackage{hyperref}
\hypersetup{colorlinks=true,linkcolor=black,urlcolor=black,pdfborder={0 0 0}}

% --- Métadonnées du lot -----------------------------------------------------
% Reprises telles quelles par le script de rendu, qui les affiche en tête de
% la vue de lecture. Elles fixent ce que le fichier prétend couvrir : sans
% elles, une transcription flotte sans point d'attache dans un fonds de seize
% mille feuillets.

\newcommand{\folder}[1]{\gdef\grothFolder{#1}}
\newcommand{\batch}[1]{\gdef\grothBatch{#1}}
\newcommand{\leaves}[2]{\gdef\grothFirst{#1}\gdef\grothLast{#2}}
\newcommand{\foldertitle}[1]{\gdef\grothFoldertitle{#1}}
\gdef\grothFolder{?}\gdef\grothBatch{?}\gdef\grothFirst{?}\gdef\grothLast{?}\gdef\grothFoldertitle{}

% --- Le résumé --------------------------------------------------------------
%
% Il ouvre la lecture modernisée, sous le titre « Résumé », et il est composé
% en petit. Ce n'est pas un ornement : c'est le seul passage du document qui
% ne soit pas des mathématiques, et un lecteur doit pouvoir le reconnaître
% avant de l'avoir lu — puis le sauter s'il connaît déjà le sujet, ou s'y
% arrêter s'il ne le connaît pas. Le filet qui le suit marque la reprise du
% corps.

\newenvironment{resume}
  {\small\setlength{\parskip}{0.45em}}
  {\par\medskip\hrule\medskip}

% --- Mots-clés ---------------------------------------------------------------
%
% En anglais, en clôture du résumé de la lecture modernisée : le vocabulaire
% moderne sous lequel chercher ce que les feuillets construisent. Le manifeste
% du site (`npm run manifest`) les extrait pour étiqueter la cote entière —
% cette ligne est la source unique des tags, il n'y a pas de fichier à côté.
\newcommand{\keywords}[1]{\par\smallskip\noindent
  {\footnotesize\textsc{Keywords} --- \textit{#1}}\par}

% --- L'appareil critique ----------------------------------------------------

% Le numéro de feuillet, en marge. C'est l'ancre qui permet de régler un
% désaccord sur une formule en regardant *un* feuillet plutôt qu'une chemise
% de six cents. Le site s'en sert aussi pour tourner les pages du fac-similé
% quand on fait défiler la transcription.
\newcommand{\leaf}[1]{%
  \marginnote{\footnotesize\color{gray!70}\textsf{#1}}%
  \label{leaf:#1}\ignorespaces}

% Illisible : on ne devine pas. Un mot inventé qui se lit comme les autres est
% le pire résultat possible.
\newcommand{\ill}{\textcolor{red!60!black}{[\,\ldots\,]}}

% Lecture douteuse : le mot est donné, et signalé comme incertain.
\newcommand{\uncertain}[1]{\uline{#1}}

% Ajout de l'éditeur — une lettre manquante, un mot sous-entendu.
\newcommand{\add}[1]{\textcolor{blue!50!black}{[#1]}}

% Biffé par Grothendieck lui-même. Ce qu'il a rayé fait partie du document :
% une rature dit souvent où la pensée a changé de direction.
\newcommand{\struck}[1]{\sout{#1}}

% Note du transcripteur, sur l'état matériel du feuillet.
\newcommand{\note}[1]{\par\smallskip{\small\color{gray!60!black}\textit{[#1]}}\par\smallskip}

% Note marginale de Grothendieck, distincte du corps du texte.
\newcommand{\marginal}[1]{\par\smallskip\noindent\hspace{1em}%
  {\small\color{gray!60!black}#1}\par\smallskip}

% --- Titre ------------------------------------------------------------------

\AtBeginDocument{%
  \begin{center}
    {\large\bfseries Fonds Alexandre Grothendieck --- cote n\textsuperscript{o}~\grothFolder}\\[2pt]
    {\itshape\grothFoldertitle}\\[4pt]
    {\small feuillets \grothFirst--\grothLast\ (lot \grothBatch)}\\[2pt]
    {\footnotesize\color{gray!60!black}Transcription établie d'après le fac-similé numérisé
     par l'Université de Montpellier.\\
     Les lectures incertaines sont soulignées, les illisibles notées [\,\ldots\,],
     les ajouts entre crochets.}
  \end{center}
  \vspace{1em}}

\endinput


\folder{16}
\batch{1}
\pages{1}{20}
\dating{1966-1967}
\watermark{Édition de démonstration}
\foldertitle{Formalisme algébrique des correspondances et algèbres de Poincaré-Hodge (cf conjectures standard) : notes manuscrites (s.d.), tapuscrit annoté (s.d.), lettres (1966-1967)}

\begin{document}

\section{Sur certains anneaux gradués}

\note{Le titre est celui qu'il inscrit, souligné, en tête du feuillet 2, qui
porte aussi sa pagination « 1 ». La suite court sans interruption jusqu'au
feuillet 17, paginée de sa main 1 à 15, avec un « 5 bis » au feuillet 7 ; sa
pagination est rappelée page par page, parce que ses renvois internes s'y
réfèrent.}

\page{2}\note{sa page 1}

Soit $k$ anneau (commut.) fixé.

1. Soient \struck{$k$ un anneau,} $M^{*}$ un $k$-module gradué (à graduation
\ill{}), $\mathcal{E}^{*}$ un \add{sous-anneau de l'anneau} \struck{une algèbre
graduée} d'endomorphismes de $M^{*}$. On désigne par $\pi_{i}$ la projection
canonique de $M^{*}$ sur $M^{i}$. \struck{et} On suppose que l'on a :

\[ (1.\ill{}) \qquad \pi_{i} \in \mathcal{E}^{*} \quad \text{pour tout } i . \]

\note{L'étiquette de cette formule est réécrite par-dessus une autre et
n'est plus lisible ; le second chiffre pourrait être un $2$. Les étiquettes
(1.1), (1.2), (1.3), (1.4) sont, elles, sûres.}

Cela implique que $\mathcal{E}^{*}$ est un sous-\struck{algèbre}
\add{anneau} gradué\supplied{e} de l'algèbre graduée \struck{des}
$\operatorname{End}(M^{*})$ (dont la graduation est définie, si l'on veut, par
\[ \operatorname{End}^{(\nu)}(M) = \textstyle\sum \operatorname{Hom}(M^{i}, M^{i+\nu}) \]
ou de façon équivalente
\[ \operatorname{End}^{(\nu)}(M) = \textstyle\sum_{i} \pi_{i+\nu} \operatorname{End}(M^{*})\, \pi_{i} , \]
la composante homogène de degré $\nu$ d'un endomorphisme $u$ de $M^{*}$ étant
\[ \struck{\mathrm{gr}}\ \mathrm{pr}_{\nu}(u) = \textstyle\sum_{i} \pi_{i+\nu}\, u\, \pi_{i} . \]

La graduation de $\mathcal{E}^{*}$ est définie (intrinsèquement, en termes de la
donnée des $\pi_{i}$), par

\[ (1.3) \qquad \text{a)}\ \ \mathcal{E}^{(\nu)} = \sum_{i} \pi_{i+\nu}\, \mathcal{E}^{*}\, \pi_{i} , \qquad \text{b)}\ \ \mathrm{pr}_{\nu}\, u = \sum_{i} \pi_{i+\nu}\, u\, \pi_{i} , \]

et

\[ (1.4) \qquad u \in \mathcal{E}^{(\nu)} \iff u\, \pi_{i} = \pi_{i+\nu}\, u \quad \text{pour tout } i . \]

\marginal{Noter les relations}

\[ (1.1) \qquad \begin{cases} \pi_{i}^{2} = \pi_{i} \\ \pi_{i} \pi_{j} = 0 \ \text{ si } i \neq j \\ \sum_{i} \pi_{i} = 1 \end{cases} \]

2. Soit $\mathcal{E}^{*}$ un \struck{$k$-algèbre} anneau gradué, et soient
$\pi_{i}$ $(i \in \mathbb{Z})$ \add{nuls sauf un nombre fini} des éléments de
$\mathcal{E}$. À quelles conditions les formules (1.3) définissent-elles

\page{3}\note{sa page 2}

sur $\mathcal{E}^{(\nu)}$ une structure d'\struck{algèbre} anneau gradué ?

Il suffit que l'on ait les relations (1.1.)

\[ \begin{cases} \pi_{i}^{2} = \pi_{i} \\ \pi_{i} \pi_{j} = \pi_{j} \pi_{i} = 0 \ \text{ si } i \neq j \\ \sum_{i} \pi_{i} = 1 \end{cases} \]

qui impliquent que les $\mathrm{pr}_{\nu}^{\mathcal{E}}$ définis par (1.2.) sont
des projecteurs deux à deux orthogonaux de somme $1$, donc que les
$\mathcal{E}^{(\nu)}$ définissent bien une graduation \add{de $k$-modules} sur
$\mathcal{E}^{*}$, \struck{puis} et que
$\mathcal{E}^{(\nu)} \mathcal{E}^{(\nu')} \subset \mathcal{E}^{(\nu + \nu')}$.

\note{Le renvoi « (1.2.) » est écrit ainsi ; la formule qui définit
$\mathrm{pr}_{\nu}$ porte, page 2, l'étiquette (1.3) b). La discordance est sur
le feuillet.}

Noter aussi que (1.4) sont valables, \struck{en \ill{}} avec \ill{}
\struck{où $u = \sum_{i,j} \pi_{j} u \pi_{i}$} \struck{de la relation
$\pi_{i+\nu} u \pi_{i} = \pi_{i} u \pi_{i}$}

\[ \begin{cases} \mathcal{E} = \bigoplus_{i,j} \pi_{j}\, \mathcal{E}\, \pi_{i} \\ \mathrm{pr}_{i,j}(u) = \pi_{j}\, u\, \pi_{i} \end{cases} \qquad \text{où } u = \sum_{\alpha,\beta} \pi_{\beta}\, u\, \pi_{\alpha} \]

et la relation $u \pi_{i} = \pi_{i+\nu} u$ implique, puisque
$u \pi_{i} = \sum_{\beta} \pi_{\beta} u \pi_{i}$ et
$\pi_{i+\nu} u = \sum_{\alpha} \pi_{i+\nu} u \pi_{\alpha}$, que
\[ \pi_{\beta}\, u\, \pi_{i} = 0 \quad \text{si } \beta \neq i + \nu . \]

\marginal{Regarder les modules sur \struck{des} l'anneau $\mathcal{E}$ muni des éléments $\pi_{i}$ comme donnés.}

3. \struck{Nous pouvons} Partons enfin d'un \struck{algèbre} anneau gradué
$\mathcal{E}^{*}$ \struck{sur $k$}, et posons-nous la question de l'existence et
de l'unicité d'une famille $(\pi_{i})_{i \in \mathbb{Z}}$ \add{satisfaisant
(1.1.)} d'\ill{} de $\mathcal{E}^{*}$, définissant la graduation, par (1.3.)
(ou par (1.4)). On suppose ici

\page{4}\note{sa page 3}

\struck{\ill{}} et \struck{que} l'on suppose que

\[ (3.0) \qquad \pi_{i} = 0 \ \text{ si } i \notin [0, 2n], \qquad n \ \text{entier \add{fixé}} . \]

\struck{Supp} On s'occupe ici de l'unicité. \struck{Supposons établi
l'unicité} Notons d'abord que
\[ u \in \mathcal{E}^{0} \iff \bigl( u \pi_{i} = \pi_{i} u \ \ \forall i \bigr) , \]
donc $\mathcal{E}^{0}$ est le commutant de la $k$-algèbre engendrée par les
$\pi_{i}$. Donc les $\pi_{i}$ sont dans le centre
$\mathfrak{z}\mathcal{E}^{0}$ de $\mathcal{E}^{0}$, et ce sont les projecteurs
\add{mutuellement} orthogonaux de somme $1$ de cet anneau commutatif.

Supposons \struck{déterminés} déjà les $\pi_{\alpha}$ pour
$\alpha \notin [i, 2n-i]$ \add{$i < n$}, et déterminons $\pi_{i}$ et
$\pi_{2n-i}$. Posons
\[ \varphi_{i} = 1 - \sum_{\alpha < i} \pi_{\alpha} = \sum_{\alpha \geqslant i} \pi_{\alpha} , \qquad \psi_{2n-i} = 1 - \sum_{\alpha > 2n-i} \pi_{\alpha} = \sum_{\alpha \leqslant 2n-i} \pi_{\alpha} , \]
donc aussi les composés
\[ \psi_{2n-i}\, \mathcal{E}^{2n-2i}\, \varphi_{i} = \sum_{\alpha} \psi_{2n-i}\, \pi_{\alpha + (2n-2i)}\, \mathcal{E}\, \pi_{\alpha}\, \varphi_{i} . \]

\marginal{C'est clair si $i = n-1$, car on connaît déjà les $\pi_{\alpha}$ pour
$\alpha \neq n$, d'où aussi $\pi_{n} = 1 - \sum_{\alpha \neq n} \pi_{\alpha}$.
Mais pour $i \leqslant n-1$ \ill{}}

Or
\[ \pi_{\alpha}\, \varphi_{i} = \begin{cases} 0 & \text{si } \alpha < i \\ \pi_{\alpha} & \text{si } \alpha \geqslant i \end{cases} \qquad \text{et} \qquad \psi_{2n-i}\, \pi_{\alpha+(2n-2i)} = \begin{cases} 0 & \text{si } \alpha > i \\ \pi_{\alpha+(2n-2i)} & \text{si } \alpha \leqslant i \end{cases} \]
donc le terme d'indice $\alpha$ dans la somme est nul sauf si $\alpha = i$,
auquel cas il est égal à
\[ \pi_{i+(2n-2i)}\, \mathcal{E}\, \pi_{i} = \pi_{2n-i}\, \mathcal{E}\, \pi_{i} = \mathcal{E}_{2n-i,\,i} . \]

Donc la connaissance des $\pi_{\alpha}$ pour $\alpha \notin [i, 2n-i]$ implique
la connaissance de $\mathcal{E}_{2n-i,\,i}$ \struck{considéré}. La condition qui
va impliquer l'unicité de $u_{i}$,

\page{5}\note{sa page 4}

$u_{2n-i}$ est \struck{sous} \ill{} la conjonction des suivantes :

\[ (3.1)_{i} \qquad \bigl( u \in \mathfrak{z}\mathcal{E}^{0}, \ \mathcal{E}^{2n-i}_{i}\, u = 0 \bigr) \implies \pi_{i}\, u = 0 \]
\[ (3.2)_{i} \qquad \bigl( u \in \mathfrak{z}\mathcal{E}^{0}, \ u\, \mathcal{E}^{2n-i}_{i} = 0 \bigr) \implies \pi_{2n-i}\, u = 0 \]

En effet, la relation $(3.1)$ implique que \struck{$\pi_{i} \in \mathfrak{z}\mathcal{E}^{0}$
est} l'annulateur de $\pi_{i}$ dans $\mathfrak{z}\mathcal{E}^{0}$ est l'idéal
des $u \in \mathfrak{z}\mathcal{E}^{0}$ tels que $\mathcal{E}^{2n-i}_{i} u = 0$ ;
et dans un anneau commutatif, un idempotent est connu quand on connaît son
annulateur. De même pour $\pi_{2n-i}$, son annulateur dans
$\mathfrak{z}\mathcal{E}^{0}$ est formé des $u \in \mathfrak{z}\mathcal{E}^{0}$
tels que $u\, \mathcal{E}^{2n-i}_{i} = 0$.

En même temps, les conditions $(3.1)_{i}$, $(3.2)_{i}$, pour tout $i$, donnent
une \struck{condition} possibilité de déterminer de proche en proche les
$\pi_{i}$.

4. Lorsque $\mathcal{E}^{*}$ est comme dans 1. \add{ou v. (2.2.)}, une condition
suffisante pour avoir $(3.1)_{i}$ est la \struck{\ill{}} suivante :

\[ (4.1)_{i} \qquad \text{l'intersection des noyaux des } u \colon M^{i} \to M^{2n-i}, \ u \in \mathcal{E}^{2n-i}_{i}, \ \text{est réduite à } 0 . \]
\[ (4.2)_{i} \qquad \text{la somme des images des } u \colon M^{i} \to M^{2n-i}, \ u \in \mathcal{E}^{2n-i}_{i}, \ \text{est égale à } M^{2n-i} . \]

Ces deux conditions s'expriment \ill{}

\[ (4.3)_{i} \qquad \text{il existe } u \in \mathcal{E}^{2n-i}_{i} \ \text{qui soit un isomorphisme} \ \ill{}, \ \text{ou, de} \]

\page{6}\note{sa page 5}

façon équivalente, il existe $u \in \mathcal{E}^{(2n-2i)}$, tel que le morphisme
induit $u \colon M^{i} \to M^{2n-i}$ soit un isomorphisme.

On peut transformer ceci de la façon suivante :

\[ (4.4)_{i} \qquad \begin{cases} \exists\, v_{i} \in \mathcal{E}^{(2n-2i)},\ w_{i} \in \mathcal{E}^{-(2n-2i)} \ \text{tels que l'on ait} \\[2pt] (w_{i} v_{i} - 1)\, \pi_{i} = 0 , \qquad (v_{i} w_{i} - 1)\, \pi_{2n-i} = 0 \end{cases} \]

Or, sans plus supposer que $\mathcal{E}^{*}$ se réalise à l'aide d'un $M^{*}$,
on voit aussitôt que $(4.4)_{i}$ implique $(3.1)_{i}$ et $(3.2)_{i}$. En effet,
prenant $v$, $w$ comme \add{ci-dessus} : $v \pi_{i} = \pi_{2n-i} v$,
$w \pi_{2n-i} = \pi_{i} w$, \struck{on voit} \struck{supposons} $v$, $w$ dans
$\mathcal{E}^{2n-i}_{i}$, $\mathcal{E}^{i}_{2n-i}$. Si alors
\struck{$u \in \mathcal{E}^{(0)}$} $u \in \mathcal{E}^{(0)}$ est tel que
$\mathcal{E}^{2n-i}_{i} u = 0$, donc $v u = 0$, on aura
$w v u \pi_{i} = 0$, i.e.\ $w v \pi_{i} u = 0$, donc $\pi_{i} u = 0$ (car
$w v \pi_{i} = \pi_{i}$). De même pour $(3.2)_{i}$, en utilisant la dernière
relation $(4.4)_{i}$.

\marginal{\struck{et 5.} v. p. 5 bis}

6. Nous allons étudier de façon plus détaillée le cas où les conditions (4.4)
sont satisfaites. Pour ceci, \struck{considérons} $\mathcal{E}^{*}$ \add{étant
un sous-anneau gradué de $\operatorname{End}(M^{*})$} \ill{} $M^{*}$,
\struck{et supposons} qu'on ait un

\[ (6.0) \qquad L \in \mathcal{E}^{(2)} \]

qui ait la propriété que \ill{}

\page{7}\note{sa page « 5 bis » ; le paragraphe 5 est écrit après coup et inséré ici}

5. On peut aussi donner une condition d'existence des $\pi_{i}$, liée aux
conditions (4.4). Pour ceci, on suppose que $\mathcal{E}$ est contenu dans un
sur-anneau gradué $\mathcal{F}$ \add{p.\ ex.\ \struck{des endom.}\ l'anneau des
endom.\ d'un $M^{*}$ comme dans 1.} pour lequel les $\pi_{i}$ satisfont (1.1.)
et les relations correspondantes (1.3) (avec $\mathcal{E}$ remplacé par
$\mathcal{F}$) \add{et $(3.0)$ : $\pi_{i} = 0$ si $i \notin [0,2n]$} sont
satisfaites. \struck{Supposons} Supposons que pour tout $i$ la condition
$(4.4)_{i}$ soit satisfaite (avec $v_{i}$, $w_{i} \in \mathcal{E}$ bien sûr).
On en conclura que

\[ (5.1) \qquad \pi_{i} \in \mathcal{E} \quad \text{pour tout } i . \]

\struck{Raisonnons} \struck{Supposons provisoirement} \ill{} $\pi_{\alpha} \in
\mathcal{E}$ si $\alpha \notin [i, 2n-i]$, \ill{} vérifier \ill{}
\struck{si $i \leqslant n$,} i.e.\ $\pi_{i}$ et $\pi_{2n-i} \in \mathcal{E}$.
(Trivial si $i = n-1$.) \struck{Supposons pour $i \leqslant n-2$}
\add{comme $\varphi_{i}$, $\psi_{2n-i} \in \mathcal{E}$, et} pour $i+1$,

Notons que, avec les notations de 3., \struck{\ill{}}
$\psi_{2n-i} v_{i} \varphi_{i}$, $\varphi_{i} w_{i} \psi_{2n-i} \in
\mathcal{E}$, où
\[ v'_{i} = \psi_{2n-i}\, v_{i}\, \varphi_{i} = \pi_{2n-i}\, v_{i}\, \pi_{i} \in \mathcal{E} , \]
\[ w'_{i} = \varphi_{i}\, w_{i}\, \psi_{2n-i} = \pi_{i}\, w_{i}\, \pi_{2n-i} \in \mathcal{E} , \]
et \struck{comme}
\[ w'_{i} v'_{i} = \pi_{i}\, w_{i}\, \pi_{2n-i}\, v_{i}\, \pi_{i} = \pi_{i}\, \pi_{i}\, w_{i}\, v_{i}\, \pi_{i} = \pi_{i}^{3} = \pi_{i} \]
\marginal{car $w_{i} \pi_{2n-i} = \pi_{i} w_{i}$ ; car $w_{i} v_{i} \pi_{i} = \pi_{i}$}
et de même $v'_{i} w'_{i} = \pi_{2n-i}$, on conclut que $\pi_{i} \in
\mathcal{E}$, $\pi_{2n-i} \in \mathcal{E}$.

\page{8}\note{sa page 6 ; la rédaction du n° 6 reprend ici}

\[ (6.1)_{i} \qquad L^{n-i} \,|\, M^{i} \colon M^{i} \to M^{2n-i} \ \text{ est un isom.} \]

Ceci implique déjà $(4.1)_{i}$ et $(4.2)_{i}$ et $(3.1)_{i}$ et $(3.2)_{i}$,
donc l'unicité des $\pi_{i}$. Supposons de plus

\[ (6.2)_{i} \qquad \ill{} \ \text{est élément de } \mathcal{E}^{-(2n-2i)} . \]

\note{Le libellé de $(6.2)_{i}$ est tracé très vite et presque entièrement
perdu sous une surcharge ; ce qui se lit est qu'il s'agit de l'inverse de
l'isomorphisme de $(6.1)_{i}$, exigé dans $\mathcal{E}$.}

\marginal{en termes des conditions (6.1) et (6.2) : $L \in \mathcal{E}^{2}$
soit un élément de \ill{} $\mathcal{E}^{*}$ \ill{} de $\mathcal{E}$, déjà les
$\pi_{i} \in \mathcal{E}^{*}$}

\[ (6.3)_{i} \qquad \begin{cases} \exists\, w_{i} \in \mathcal{E}^{-(2n-2i)} \ \text{tel qu'on ait} \\[2pt] (w_{i} L^{n-i} - 1)\, \pi_{i} = 0 , \qquad (L^{n-i} w_{i} - 1)\, \pi_{2n-i} = 0 \end{cases} \]

Cette condition \ill{} : elle est indépendante de la donnée particulière d'un
$M^{*}$. \struck{\ill{} donc (3.1) et (3.2)} Elle \ill{} $(4.4)_{i}$, et
(comme on a vu aux n° 4, 5) implique $\pi_{i} \in \mathcal{E}$ et l'unicité des
$\pi_{i}$, en termes des conditions (4.1) et (4.2).

Revenons à $M^{*}$, et $L$ satisfaisant (6.1). Nous avons \ill{} comment
\ill{} cette condition implique la décomposition primitive de $M^{*}$ :

\[ (6.4) \qquad \begin{cases} M^{i} \simeq P^{i} + L P^{i-2} + L^{2} P^{i-4} + \ldots & (i \leqslant n) \\[2pt] M^{2n-i} \simeq L^{n-i} P^{i} + L^{n-i+1} P^{i-2} + \ldots \end{cases} \]

On désigne par $\pi_{i,\alpha}$ la projection de $M^{i}$ sur
$L^{\alpha} P^{i-2\alpha}$, par $\pi_{2n-i,\alpha}$ celle de $M^{2n-i}$ sur
$L^{n-i+\alpha} P^{i-2\alpha}$.

\page{9}\note{sa page 7}

\struck{(6.5) \ill{}}

\textbf{Théorème.} Les conditions suivantes sont équivalentes \add{(i) à (iii)} :

\begin{enumerate}
  \item Les conditions $(6.1)_{i}$ et $(6.2)_{i}$ pour tout $i$, et les $\pi_{i} \in \mathcal{E}$.
  \item \struck{\ill{}} donc la condition $(6.3)_{i}$ pour tout $i$.
  \item \ill{} $\Lambda \in \mathcal{E}^{(-2)}$ \ill{}
  \item Les $\pi_{i,\alpha}$ sont dans $\mathcal{E}$ \ill{}
  \item Les $\pi_{i} \in \mathcal{E}$, \ill{} et (1.3) a) et b), ou \ill{} (1.4).
  \item L'élément $\Lambda$ de $\mathcal{E}^{(-2)}$ est déterminé de façon unique
        \struck{\ill{}} (en termes de $\mathcal{E}^{*}$ et de $L \in \mathcal{E}^{(2)}$),
        par les conditions suivantes :
\end{enumerate}

\note{La liste est très raturée : les énoncés (ii) à (v) sont écrits, biffés,
réécrits et renumérotés en marge, et seule leur charpente se lit. Ce qui est
sûr est que la chaîne va de la condition de Lefschetz (6.1) à l'existence dans
$\mathcal{E}$ de l'opérateur $\Lambda$ de degré $-2$, puis des projecteurs
$\pi_{i,\alpha}$ de la décomposition primitive.}

\[ (6.7) \qquad \begin{cases} L \Lambda = 1 - \sum_{0 \leqslant i \leqslant n} \pi_{i,0} \\[2pt] \Lambda L = 1 - \sum_{n \leqslant i \leqslant 2n} \pi_{i,0} \end{cases} \]

\emph{Dém.} L'équivalence de (i) et (iii) est claire, et le fait que
(i) \ill{} (iii) (en prenant $w_{i} = \Lambda^{2n-i}$). \ill{}

\page{10}\note{sa page 8}

Prouvons (iv). Supposons prouvé, pour $j \notin [i, 2n-i]$, que

\[ (6.7) \qquad \pi_{j,\alpha} \in \mathcal{E} \quad \text{pour } j \notin [i, 2n-i] \]

où $i \leqslant n$, et prouvons que pour $j = i$ ou $j = 2n-i$

\[ (6.7') \qquad \pi_{j,\alpha} \in \mathcal{E} . \]

\note{Les deux étiquettes (6.7) se suivent : celle-ci porte le même numéro que
celle de la page 9. Il ne corrige pas.}

\ill{} On raisonne par récurrence descendante sur $\alpha$ : supposons

\[ (6.8) \qquad \pi_{i,\alpha'} \in \mathcal{E}, \quad \pi_{2n-i,\alpha'} \in \mathcal{E} \quad \text{si } \alpha' > \alpha \]

et on en conclura

\[ (6.8') \qquad \pi_{i,\alpha} \in \mathcal{E}, \quad \pi_{2n-i,\alpha} \in \mathcal{E} . \]

Ceci \struck{\ill{}} prouvera bien $(6.7')$, puisqu'on a $\pi_{i,\alpha} =
\pi_{2n-i,\alpha} = 0$ pour $\alpha$ \ill{}, donc, par la relation

\[ (6.9) \qquad \sum_{\alpha} \pi_{i,\alpha} = \pi_{i} , \]

$\pi_{i,0} = \pi_{i} - \sum_{\alpha \geqslant 1} \pi_{i,\alpha}$. \ill{}

Posons
\[ q_{i,\alpha} = 1 - \sum_{\alpha' > \alpha} \pi_{i,\alpha'} = \sum_{\alpha' \leqslant \alpha} \pi_{i,\alpha'} , \]
donc $q_{i,\alpha} \in \mathcal{E}$ par (6.8).

\page{11}\note{sa page 9}

Alors on veut que
\[ L^{n-i+\alpha}\, q_{i,\alpha} \colon M^{i} \longrightarrow M^{2n-i+2\alpha} \]
soit nul sur les $L^{\alpha'} P^{i-2\alpha'}$ pour $\alpha' \neq \alpha$, et que
sur $L^{\alpha} P^{i-2\alpha}$ il induise \ill{} le composé avec l'isomorphisme
\[ L^{n-i+2\alpha} \colon P^{i-2\alpha} \longrightarrow L^{n-i+2\alpha} P^{i-2\alpha} \subset M^{2n-i+2\alpha} . \]
\ill{} il s'ensuit aussitôt \ill{} que

\[ (6.9) \qquad \pi_{i,\alpha} = L^{\alpha}\, w_{i-2\alpha}\, L^{n-i+\alpha}\, q_{i,\alpha} = L^{\alpha}\, w_{i-2\alpha}\, L^{n-i+\alpha}\Bigl( 1 - \sum_{\alpha' > \alpha} \pi_{i,\alpha'} \Bigr) \]

ce qui donne bien $\pi_{i,\alpha} \in \mathcal{E}$. \struck{Ceci prouve (iv).}
Si $i < n$, il faut de même $\pi_{2n-i,\alpha} \in \mathcal{E}$. \ill{} On a
$\pi_{2n-i,\alpha} \colon M^{2n-i} \to M^{2n-i+2\alpha}$ \ill{} de la façon

\[ (6.10) \qquad \pi_{2n-i,\alpha} = L^{n-i}\, \pi_{i,\alpha}\, w_{i} \ \struck{\ill{}} \]

\emph{Reste} : prouver (iii) $\Lambda \in \mathcal{E}$. Comme on a

\[ (6.11) \qquad \Lambda = \sum_{j} \struck{\ill{}}\ \Lambda\, \pi_{j,\alpha}\ \struck{\ill{}} \]

il suffit de prouver que $\Lambda \pi_{i,\alpha} \struck{\ill{}} \in \mathcal{E}$.

\page{12}\note{sa page 10}

Or si $i \leqslant n$, on aura

\[ (6.12) \qquad \begin{cases} \Lambda\, \pi_{i,0} = 0 , \\[2pt] \Lambda\, \pi_{i,\alpha} = \struck{L}\ w_{i-2\alpha}\, L^{n-i+\alpha}\, \pi_{i,\alpha} \quad \text{si } \alpha \geqslant 1 , \end{cases} \]

\struck{et si $j > n$, \ill{} (6.13) $\Lambda \pi_{j,\alpha} = L \ill{}$}

\[ (6.13) \qquad \Lambda\, \pi_{2n-i,\alpha} = L^{n-i+\alpha-1}\, w_{i-2\alpha}\, \struck{\ill{}}\ L^{\alpha}\, \pi_{2n-i,\alpha} \qquad (\text{si } i \leqslant n-1) \]

Les conditions envisagées indiquent \ill{} (iv), i.e.\ prouvons que l'on a
alors $\pi_{i} = \pi'_{i}$. (iv) indiquent (ii). Il suffit de \ill{}
\emph{Démonstration à fournir.}

\emph{Remarque.} Voici une \ill{} variante de \ill{} unique \ill{} en termes
de $L$, et \ill{} une caractérisation \ill{} en termes de l'anneau gradué
$\mathcal{E}$) :

\marginal{$\Lambda \in \mathcal{E}^{(2)}$, $\pi_{i,\alpha} \in \mathcal{E}^{(0)}$}

\note{En marge il écrit $\mathcal{E}^{(2)}$ pour $\Lambda$, alors que partout
ailleurs $\Lambda$ est de degré $-2$ ; le signe manque sur le feuillet.}

\begin{enumerate}
  \item $\bigl( \Lambda^{n-i} L^{n-i} - 1 \bigr)\, \pi_{i} = 0$ ;
  \item $\bigl( L^{n-i} \Lambda^{n-i} - 1 \bigr)\, \pi_{2n-i} = 0$ ;
  \item $\pi_{i,\alpha} = L^{\alpha}\, \Lambda^{n-i+2\alpha}\, L^{n-i+\alpha} \Bigl( 1 - \sum_{\alpha' > \alpha} \pi_{i,\alpha'} \Bigr)$ ;
  \item $\pi_{2n-i,\alpha} = L^{n-i}\, \pi_{i,\alpha}\, \Lambda^{n-i}$ .
\end{enumerate}

\page{13}\note{sa page 11}

\[ (6.15) \qquad \begin{cases} L \Lambda = 1 - \sum_{0 \leqslant i \leqslant n} \pi_{i,0} \\[2pt] \Lambda L = 1 - \sum_{n \leqslant j \leqslant 2n} \pi_{j,0} \end{cases} \]

\struck{Notant que \ill{} $\pi_{j,0}$, il y \ill{} déterminés par les \ill{}}

\marginal{il suffit \ill{} $j = n$ \ill{} $\Phi$ : explicitée \ill{} relations \ill{}}

\emph{Question.} Donner une \ill{} synthétique des relations entre les
$\pi_{i}$ comme \struck{\ill{}} $\Lambda$, $L$ exclusivement, et \ill{} des
relations analogues \ill{} $L$, $\Lambda$ \ill{} qui définissent $\Lambda$,
polynomiales en $L$, $\Lambda$ \ill{} $\pi_{i}$ \ill{} qui donne absolument
$\mathcal{E}$.

7. Revenons à \ill{} un \ill{} $L \in \mathcal{E}^{(2)}$. \ill{} On veut
expliciter le \ill{} \ill{} utilisant $M^{*}$, \ill{} un anneau gradué
$\mathcal{E}^{*}$. Partons pour ceci \ill{} de degré $0$ à $2n$ sur
$\mathbb{Z}$ \add{sur $\mathbb{Z}$ !}, \ill{} une digression \ill{}
$\xi_{a} \in \Lambda^{2} M_{a}$, \ill{} $M^{\alpha}_{0}$, et considérons \ill{}
$\xi^{*} \in \Lambda^{2} M^{\vee}_{a}$, la forme \ill{} dont
$M^{*}_{0} = \Lambda M_{a}$, \ill{} $M^{\alpha}_{0}$ \ill{}

\note{Le début du n° 7 est écrit très vite et sa syntaxe est perdue ; ce qui
s'en dégage est qu'il part d'un module libre $M_{0}$ de rang fini sur
$\mathbb{Z}$, d'une forme $\xi_{0} \in \Lambda^{2} M_{0}$ et de sa duale
$\xi^{*}_{0}$, et prend pour $M^{*}_{0}$ l'algèbre extérieure $\Lambda M_{0}$.}

et considérons sur $M^{*}_{0}$ les opérateurs
\[ L_{0}\, x = \xi_{0}\, x , \qquad \Lambda_{0}\, x = \xi^{*}_{0} \lrcorner\, x \]
et considérons l'anneau \add{gradué}
\[ \Phi_{0} = \mathbb{Z}[L_{0}, \Lambda_{0}] \subset \operatorname{End}(M^{*}_{0}) . \]
Notons que $\operatorname{End}_{\mathbb{Z}}(M^{*}_{0})$ \add{opérant sur
$M^{*}_{\mathbb{Z}}$}, \ill{} \uncertain{muni} de $L$, satisfait à la condition
de $(6.1)$\uncertain{$(\omega)$}, \struck{\ill{}} et que $\Lambda_{0}$ est
l'opérateur associé défini par (6.15).

\note{Le $\mathbb{Z}$ est porté sous $\operatorname{End}$, et « opérant sur
$M^{*}_{\mathbb{Z}}$ » est une insertion au-dessus de la ligne, rattachée par un
long trait. Ce que la parenthèse après $(6.1)$ contient est un seul glyphe
arrondi, lu $\omega$ et non identifié autrement. La phrase identifie les deux
$\Lambda_{0}$ du feuillet : celui que la formule affichée définit par la
contraction $\xi^{*}_{0} \lrcorner\,-$, et celui que caractérisent les relations
(6.15).}

\page{14}\note{sa page 12}

\ill{} $\Phi_{0}$ \ill{} est engendré par les $\varphi^{\beta,\alpha}_{i}$
\struck{\ill{}} \struck{$L^{\beta} \pi_{i,\alpha}$}
\struck{$\varphi^{\beta}_{i,\alpha}$ \ill{}}, avec

\[ 0 \leqslant i \leqslant n , \quad 0 \leqslant \alpha, \beta \leqslant 2n-i \]

\note{À droite, biffé d'un trait, un premier jeu d'indices : « $0 \leqslant i
\leqslant 2n$, $0 \leqslant j \leqslant 2n$, $\alpha, \beta \geqslant 0$,
$i - 2\alpha = j - 2\beta \geqslant 0$ ».}

\[ \begin{cases} \varphi^{\beta,\alpha}_{i} = L_{0}^{\beta-\alpha}\, \pi_{i+2\alpha,\alpha} & \text{si } \alpha \leqslant \beta \\[2pt] \varphi^{\beta,\alpha}_{i} = \Lambda_{0}^{\alpha-\beta}\, \pi_{i+2\alpha,\alpha} & \text{si } \alpha \geqslant \beta \end{cases} \]
\[ \begin{cases} = L_{0}^{\beta-\alpha}\, \pi_{i,\alpha} = \pi_{j,\beta}\, L_{0}^{\beta-\alpha} & \text{si } \alpha \leqslant \beta \\[2pt] = \Lambda_{0}^{\alpha-\beta}\, \pi_{i,\alpha} = \pi_{j,\beta}\, \Lambda_{0}^{\alpha-\beta} & \text{si } \alpha \geqslant \beta \end{cases} \]

Les $\varphi^{\beta,\alpha}_{i}$ et leur multiplication \ill{} sont donnés
\ill{}. La multiplication \ill{} ne dépend pas du choix de \ill{} de
$(M_{0}, \xi_{0})$ \ill{} l'\ill{}, dans \ill{} $\Phi_{0}$ \ill{} bien
déterminé \ill{}.

\textbf{Proposition.} Revenant à la donnée de $M^{*}$ comme \ill{}
\struck{$n - 2g$} \add{anneau} gradué des \struck{\ill{}}
$\operatorname{End}(M^{*})$, \ill{} et de $L \in \mathcal{E}^{(2)}$.
Conditions équivalentes :

\begin{enumerate}
  \item \ill{} $L$ \ill{} $M^{*}$ satisfait les conditions du théorème du n° 6 ;
  \item \struck{Hodge} \ill{} $\Lambda$ \ill{} satisfaisant les conditions du théorème du n° 6 ;
  \item il existe un \ill{} homomorphisme d'anneaux gradués
        $\varphi \colon \Phi_{0} \to \mathcal{E}$
\end{enumerate}

\page{15}\note{sa page 13}

\ill{} assignant $L_{0}$ à $L$, et les $\pi_{i}$ \ill{} des $\pi_{i}$
\struck{\ill{}} \add{des $\pi_{i} \in \mathcal{E}$}. De plus, sous ces
conditions, l'homomorphisme \ill{} est unique, et avec les notations du n° 6,
il transforme $\Lambda_{0}$ en $\Lambda$, les $\pi_{i}$ en $\pi_{i}$, les
$\pi_{i,\alpha}$ en $\pi_{i,\alpha}$, les $\varphi^{\beta,\alpha}_{i}$ en
$\varphi^{\beta,\alpha}_{i}$ \struck{\ill{}}.

\marginal{condition essentielle !}

\emph{Dém.} (i) $\Rightarrow$ (ii) : \ill{} multiplication \ill{} donnée au
n° 6. Résulte du \ill{} théorème des $\varphi^{\beta,\alpha}_{i}$ donné au
n° 6, et du théorème \ill{}.

(ii) $\Rightarrow$ (i) : soient \ill{} $\pi'_{i}$, \ill{}
$L^{\alpha} \pi'_{i} = \pi'_{i+\alpha} L^{\alpha}$ \ill{} de $\mathcal{E}$
\ill{}. Je dis que $\pi'_{i} = \pi_{i}$ \ill{}. Comme $i \leqslant n$, on a
\[ \bigl( \Lambda_{0}^{n-i} L_{0}^{n-i} - 1 \bigr)\, \pi_{i} = 0 , \qquad \bigl( L_{0}^{n-i} \Lambda_{0}^{n-i} - 1 \bigr)\, \pi_{2n-i} = 0 , \]
\ill{} appliquant $\varphi$ (et pour $\Lambda' = \varphi(\Lambda)$)
\[ \bigl( \Lambda^{n-i} L^{n-i} - 1 \bigr)\, \pi_{i} = 0 , \qquad \bigl( L^{n-i} \Lambda^{n-i} - 1 \bigr)\, \pi_{2n-i} = 0 , \]
(ii) \ill{} du théorème du n° 6 \ill{} simplifie \ill{} définition (i).
\ill{} il résulte \struck{\ill{}} $\varphi(\Lambda) = \Lambda$, et bien
\struck{$\varphi(\pi_{i,\alpha}) = \pi_{i,\alpha}$}, \ill{} la forme
\ill{} simultanément \ill{} (6.14).

\page{16}\note{sa page 14}

\textbf{Cor.} Soit $\mathcal{E}$ un anneau \ill{} gradué, \add{et soit
$L \in \mathcal{E}^{(2)}$}. S'il existe un homomorphisme d'anneaux gradués
\[ u \colon \Phi_{0} \longrightarrow \mathcal{E} \]
tel que $u(L_{0}) = L$, et \ill{} $u(\pi_{i}) = \pi_{i}$ \ill{} exprimant la
graduation de $\mathcal{E}$ par (1.3), alors \ill{} un tel \ill{} est unique
et uniquement déterminé, et les $\pi_{i}$ sont déterminés \ill{} de la
structure d'anneau gradué de $\mathcal{E}$ \ill{} les $\pi_{i}$ satisfaisant
les conditions (1.1) et (1.3).

\ill{} l'existence \ill{} sont alors équivalentes : l'existence \ill{}
$\exists\, w_{i} \in \mathcal{E}^{-(2n-2i)}$ tel que \ill{} $0 \leqslant i
\leqslant n$ \ill{} et alors \ill{}
\[ \bigl( w_{i} L^{n-i} - 1 \bigr)\, \pi_{2n-i} \struck{\ill{}} = 0 , \]
\ill{} $\Lambda = u(\Lambda)$, \struck{et} $\pi_{i,\alpha} = u(\pi_{i,\alpha})$.

Comme \struck{\ill{}} des $L$ et $M^{*}$, \ill{} $\mathcal{E}^{*}$ \ill{} les
$\pi_{i}$. On convient : $\exists$ un unique \ill{} graduation \ill{}
(6.14). \ill{} $\mathcal{E}$, l'existence de la question \ill{} : l'existence
\ill{} fidèle sous $\mathcal{E}^{*}$, tel que \ill{} $\xi^{*}$ \ill{} la
donnée \ill{} les $\pi_{i} = u(\pi_{i})$, \ill{} la graduation de \ill{}.

\emph{Question.} Si \ill{} (elle sera obtenue à partir \ill{}) $\mathcal{E}$
\ill{} primitif \ill{} \ill{} contre-exemple \ill{} un hom
$\Phi_{0} \to \mathcal{E}$ tel que $u(L_{0}) = L$ ?

\marginal{Non}

déjà dans le cas $n = 1$, $2n = 2$, où il y a \ill{} $\Phi_{0}$ formé de
$\pi_{0}$, $\pi_{1}$, $\pi_{2}$, $L$, $\Lambda$, et \ill{} que
$\mathbb{Z} + \mathbb{Z}L$ ! Il faut donc \struck{en même temps} $L$ et les
$\pi_{i}$.

\page{17}\note{sa page 15}

8) Soit maintenant \struck{\ill{}} $\mathcal{E}^{*}$ un anneau gradué,
\ill{} \struck{\ill{}} et soit $L \in \mathcal{E}^{(2)}$ satisfaisant aux
\emph{conditions de Lefschetz} (de les conditions du \ill{}).

\note{Il écrit « conditions de Lefschetz » en toutes lettres et le souligne :
c'est le nom qu'il donne ici à ce que le n° 6 a caractérisé.}

\ill{} combinaison des \ill{} un anneau \ill{} correspondant \ill{}
(Cf.\ \struck{Lefschetz} \ill{}).

N.B. À montrer que \ill{} $L$, \ill{} on déduit $\Lambda$ de \ill{}, pour
\ill{} symétrique $\mathcal{E}^{*}$ \ill{}, et pour \ill{}
\struck{\ill{}} \ill{} permettant \ill{}.

\emph{Conditions.} Bornons-nous \ill{} : voici les conditions nécessaires (par
un accomplissement \ill{}) \ill{} de $L$. Une condition nécessaire est
\ill{}, i.e.

\begin{enumerate}
  \item $\mathcal{E}^{(2n)} \neq 0 \iff \pi_{0} \neq 0$ \ill{}
  \item $\pi_{2n} \neq 0$
  \item $L^{n} \neq 0$
\end{enumerate}
\[ \mathcal{E}^{(2n)} \neq 0 \iff \pi_{2n} \neq 0 \implies L^{n} \neq 0 \implies \Lambda^{n} \neq 0 . \]

\marginal{déjà pour $n$}

\note{Le feuillet s'arrête ici, au milieu de la liste des conditions
nécessaires ouverte par le n° 8 : l'argument se poursuit au-delà de la page 20
et n'est pas achevé dans ce lot.}

\section{Feuillets détachés}

\note{Les trois derniers feuillets du lot ne font pas partie du manuscrit
précédent. Les pages 18 et 20 sont deux feuillets d'un tapuscrit de sa main,
numérotés 45 et 42 dans sa propre pagination, dont le sujet est la cohomologie
$\ell$-adique, le groupe de Picard et le groupe de Néron-Severi, et qu'il a
corrigé et annoté à l'encre ; ils sont transcrits ici parce qu'ils portent ces
corrections, mais leur numérotation propre montre assez qu'ils viennent
d'ailleurs. La page 19 est le verso d'un tel feuillet et porte, de sa main, une
page de calculs sur $L$ et $\Lambda$ qui relèvent, eux, du sujet du lot. Ses
insertions interlinéaires sont données entre crochets, ses renvois marginaux en
retrait, les suppressions du tapuscrit barrées.}

\page{18}\note{page 45 du tapuscrit}

en groupes \add{lisse et} \struck{localement} de type fini sur $k$, soit
$\underline{\mathrm{Pic}}_{U}$, et un isomorphisme

\[ (8.13) \qquad \mathrm{Pic}(U) = \underline{\mathrm{Pic}}_{U}(k) \]

de $\mathrm{Pic}(U)$ avec le groupe des points à valeurs dans $k$ de
$\underline{\mathrm{Pic}}_{U}$. Cela nous fournit donc pour $\mathrm{Pic}(U)$
une structure d'extension

\[ (8.14) \qquad 0 \longrightarrow \mathrm{Pic}(U)^{0} \longrightarrow \mathrm{Pic}(U) \longrightarrow \mathrm{NS}(U) \longrightarrow 0 , \]

où on a posé

\[ (8.15) \qquad \mathrm{Pic}(U)^{0} = \underline{\mathrm{Pic}}_{U}^{0}(k) , \qquad \mathrm{NS}(U) = \underline{\mathrm{NS}}_{U}(k) , \]

$\underline{\mathrm{Pic}}_{U}^{0}$ désignant la composante neutre de
$\underline{\mathrm{Pic}}_{U}$, et $\underline{\mathrm{NS}}_{U}$ le quotient de
$\underline{\mathrm{Pic}}_{U}$ par cette composante. La construction indiquée
dans loc.\ cit.\ permet d'ailleurs de préciser la structure de $\mathrm{NS}(U)$,
en introduisant la matrice d'intersection

\[ (C_{i} . C_{j})_{1 \leqslant i, j \leqslant r} \]

définie par les $r$ composantes irréductibles $C_{i}$ de $f^{-1}(x)_{\text{réd}}$
sur le schéma régulier $X'$ : en vertu de MUMFORD \add{28}, cette matrice
est définie négative, donc le conoyau de l'homomorphisme

\[ \underline{\mathbb{Z}}^{r} \longrightarrow \underline{\mathbb{Z}}^{r} \]

qu'elle définit est un groupe fini, dont l'ordre n'est autre que la valeur
absolue du déterminant de cette matrice. Ceci posé, on trouve que
$\mathrm{NS}(U)$ est canoniquement isomorphe à ce groupe fini. La matrice
envisagée étant symétrique, on voit aussitôt que ce groupe est canoniquement
autodual (par un accouplement \struck{symétrique} à valeurs dans
$\mathbb{Q}/\mathbb{Z}$), d'où une autodualité \add{symétrique}

\[ (8.16) \qquad \mathrm{NS}(U) \times \mathrm{NS}(U) \longrightarrow \mathbb{Q}/\mathbb{Z} , \]

dont on constate aisément qu'elle est indépendante de la résolution $X'$
choisie.

Si maintenant $\ell$ est un nombre premier distinct de car $k$, alors la
structure connue des groupes algébriques commutatifs connexes sur $k$
\add{9, exp.\ 4} montre que $\mathrm{Pic}(U)^{0}$ est un groupe
$\ell$-divisible, donc

\page{19}\note{feuillet de calculs, de sa main, au dos d'un feuillet du tapuscrit}

\ill{} $b_{0}$, $b_{1}$, $b_{2}$, $b_{3}$, $b_{4}$

\[ n b_{0} + (n-1) b_{1} = \ldots , \qquad 2 b_{0} + b_{1} , \qquad \sum_{i} (n-i)\, b_{i} \]

\[ M^{0} \xrightarrow{\ L\ } M^{2n} \]

\[ L , \qquad \mathcal{E}^{(n-i)} , \qquad \Lambda , \qquad \Lambda L \]

\[ L \sim \lambda L , \qquad \Lambda = P(L \ , \qquad \Lambda = \Omega(L)\, \Delta(L) \]

\note{Ces deux dernières lignes sont écrites l'une au-dessus de l'autre et
sans membre de droite achevé : il essaie d'exprimer $\Lambda$ comme un
polynôme en $L$.}

\[ L , L' \quad | \quad L^{2}, L L' \quad | \quad L^{3} \]
\[ L'^{2} = L L' , \qquad L^{3} = L^{2} L' , \qquad L'^{3} = L' L'^{2} = L^{2} L' = L^{3} \]
\[ \bigl( L^{q},\ L'^{2} - L L',\ L^{2} L' - L^{3},\ L^{3} L' \bigr) \]
\[ \mathbb{Z}[L]/(L^{4}) \ \ill{}\ \mathbb{Z}[L'] / (L'^{2} - L'') \]

\note{Le feuillet n'est qu'une suite de calculs sans phrase liante : il essaie
sur un exemple à deux générateurs $L$, $L'$ de degré $2$ ce que les nos 6 et 7
demandent de l'anneau engendré par $L$, en écrivant une base graduée
$L, L' \mid L^{2}, LL' \mid L^{3}$ et les relations qui la ferment. La dernière
ligne, où il compare cet anneau à $\mathbb{Z}[L]/(L^{4})$, n'est pas lisible
jusqu'au bout.}

\page{20}\note{page 42 du tapuscrit}

on trouve une suite exacte canonique (où
$\underline{\mathbb{Z}}_{\ell}[1] = \varprojlim_{\nu} \mu_{\ell^{\nu}}$
est le faisceau \add{$\ell$-adique de TATE}) :

\[ (8.3) \qquad 0 \longrightarrow t^{i-1}(X, \ell) \longrightarrow H^{i}(X, \underline{\mathbb{Z}}_{\ell}[1]) \longrightarrow T_{\ell}\bigl( H^{i}(X, \mu_{\ell^{\infty}}) \bigr) \longrightarrow 0 , \]

où \struck{\ill{}} on a posé

\[ (8.4) \qquad t^{i-1}(X, \ell) = H^{i-1}(X, \mu_{\ell^{\infty}}) \big/ H^{i-1}(X, \mu_{\ell^{\infty}})^{0} \]

qui est donc un groupe fini de $\ell$-torsion, et s'identifie donc, par (8.3),
au sous-\struck{groupe} \add{module} de torsion du
$\underline{\mathbb{Z}}_{\ell}$-module de type fini

\[ (8.5) \qquad H^{i}(X, \underline{\mathbb{Z}}_{\ell}[1]) = \varprojlim H^{i}(X, \mu_{\ell^{\nu}}) , \]

dont la partie libre s'identifie, à son tour, à
$T_{\ell}\bigl( H^{i}(X, \mu_{\ell^{\infty}}) \bigr)$. On en déduit également
un isomorphisme

\end{document}
